The Soil and Water Assessment Tool (SWAT) and its restructured successor SWAT+
\citep{Arnold1998SWAT,Bieger2017SWATPlus} are among the most widely used process-based
watershed models. As input data and computing have improved, the community has pushed SWAT+
toward larger domains and hyper-resolution discretizations---tens of thousands of
hydrologic response units (HRUs) and channels in a single basin. At that scale the cost of a
\emph{single} simulation becomes the binding constraint. The benchmark model used throughout
this paper, a USGS HUC-8 watershed (hydrologic unit 03100101) is a regional, hyper-resolution model
generated by the SWATGenX platform; it contains 57{,}998 HRUs and 8{,}181 channels (full
specification in Section~\ref{sec:model}), and one simulated year takes about three minutes on
a single core. Calibration---the dominant computational activity in practice---requires
hundreds to thousands of such evaluations, so per-run cost multiplies directly into days of
wall-clock time and substantial cloud expense.

\subsection{The gap: engine-level vs ensemble-level parallelism}
Most prior work on ``parallel SWAT'' accelerates the \emph{calibration}, not the
\emph{model}. Frameworks distribute independent parameter sets across cores or cluster nodes
\citep{Rouholahnejad2012Parallelization,Zhang2013Multiobjective,Yalew2013Grid}, and more
recent efforts containerize the same ensemble pattern for reproducibility
\citep{PASS4SWAT2024}. This ensemble parallelism is essential and we use it, but it has a
hard ceiling: population-based calibration of a model with a few tens of adjustable
parameters saturates at a few tens of concurrent evaluations \citep{Rafiei2022Calibration}---beyond
that, a larger population does not improve the optimum and can degrade it. Once an individual model
is large,
the only way to apply more than that handful of cores to \emph{one} model is to parallelize
the engine itself.

Engine-internal parallelization of SWAT is comparatively rare and has targeted the
embarrassingly parallel land phase: OpenMP over HRUs in the gridded SWATG variant
\citep{Zhang2017SWATG} and compiler-assisted OpenMP in SWAT \citep{Ki2015OpenMPSWAT}. The
land phase is the easy half. The hard half---channel routing---is a strictly ordered
traversal of the river network: a reach cannot be solved until its upstream contributors
are. This data dependency is exactly what makes routing the serial bottleneck, and it is the
same structural obstacle faced by every distributed routing model
\citep{David2011RAPID,Mizukami2021mizuRoute,Yamazaki2011CaMaFlood}.

\subsection{Two non-negotiable constraints}
A parallel engine intended for real use must satisfy two constraints that the existing
literature largely sidesteps. First, results must not change: a calibrated,
published, or regulatory model cannot silently produce different numbers because it ran on
more cores. Second, correctness must be demonstrable, not asserted---data races in
shared-memory scientific codes are notoriously intermittent and escape casual testing
\citep{Atzeni2016ARCHER}. We treat both as first-class: the one-thread path reproduces the
unmodified serial engine \emph{bit for bit}, and we validate the parallel paths by
\emph{thread-count invariance}, a per-output bitwise consistency criterion in the spirit of
ensemble consistency testing for Earth-system models \citep{Baker2015CESMECT}, which here
doubles as a precise race detector.

\subsection{Contributions}
This paper makes the following contributions:
\begin{enumerate}
  \item A reproducible acceleration methodology for a large process-based Fortran model:
        profile $\rightarrow$ classify $\rightarrow$ fix-by-class $\rightarrow$ validate
        $\rightarrow$ re-measure, with each optimization an independently validated commit.
  \item A bottleneck taxonomy of SWAT+ at scale, with the central empirical finding
        that the dominant cost is \emph{overhead, not hydrology}---per-object array zeroing,
        struct copies, $\mathcal{O}(N^2)$ name lookups, and a serial routing critical path.
  \item A routing-DAG-aware shared-memory parallelization: a topological level
        wavefront over the engine command order, with reentrancy achieved by privatizing
        current-object state (\texttt{threadprivate}), and a discussion of the
        Fortran \texttt{SAVE}-locals hazard that makes na\"ive privatization unsafe.
  \item A correctness framework: byte-identity at one thread plus thread-count
        invariance as an automatable race detector, and an honest accounting of a compiler-induced
        in-stream particulate race ($15$--$20\%$ on three sediment-associated species) and its
        resolution by a byte-identical mode.
  \item An architectural-limit result: a measured characterization of where the
        SWAT+ daily step is parallelizable (the wide land phase) and where it is bound by the
        converging main-stem critical path, validated across multiple CPU microarchitectures.
\end{enumerate}

The remainder of the paper reviews related work (Section~\ref{sec:related}), details the
methodology and the parallelization (Section~\ref{sec:methods}), reports scaling and
correctness results (Section~\ref{sec:results}), and discusses the architectural limits and
portability (Sections~\ref{sec:discussion}--\ref{sec:conclusions}).
